
Supposons en particulier que @@M0@@~; alors le @@M1@@-ième polynôme universel de @@M2@@ n'a pas de composantes de poids @@M3@@. Considérons maintenant le morphisme canonique~: @@M4@@~; par composition avec @@M5@@, on obtient un morphisme~: @@M6@@, qui d'après la démonstration de 1.12 se factorise par @@M7@@ et définit un morphisme @@M8@@~; on voit aussitôt que le @@M9@@-ième polynôme universel de @@M10@@ est le même que celui de @@M11@@, c'est-à-dire la composante isobare de poids @@M12@@ du @@M13@@-ième polynôme universel de @@M14@@.

Les mêmes remarques s'appliquent au cas @@M15@@~: on peut définir un morphisme de foncteurs @@M16@@ (avec @@M17@@, @@M18@@ étant le foncteur défini sur @@M19@@ par @@M20@@) de façon analogue à 1.11; on recopie alors raisonnement et conclusions de l'alinéa précédent.

1.16. Terminons sur un exemple. On suppose que @@M21@@. On définit un endomorphisme du foncteur @@M22@@, noté @@M23@@, et appelé caractère de Todd, en posant~:
@@M24@@
Les conditions i) et ii) du théorème 1.13 se vérifient immédiatement, ce qui montre que nous avons bien défini ainsi un endomorphisme de @@M25@@.

\section*{2. Définition des @@M26@@-anneaux; exemples}

Dans ce paragraphe, l'anneau de base du paragraphe 1 sera @@M27@@ (sauf en 2.9)~; on se place sur la catégorie @@M28@@ (1.8) et on note @@M29@@ la restriction de @@M30@@ à @@M31@@ (1.10).

